\documentclass[10pt,a4paper]{article}

% --- XeLaTeX: fonts + Russian ---
\usepackage{fontspec}
\usepackage{polyglossia}
\setdefaultlanguage{russian}
\setotherlanguage{english}
\setmainfont{PT Sans}
\setmonofont[Scale=MatchLowercase]{PT Mono}
\newfontfamily\cyrillicfonttt[Scale=MatchLowercase]{PT Mono}

% --- layout ---
\usepackage[a4paper,top=0.6cm,bottom=0.55cm,left=1.3cm,right=1.3cm]{geometry}
\usepackage{titlesec}
\usepackage{enumitem}
\usepackage{xcolor}
\usepackage[hidelinks]{hyperref}

\definecolor{rule}{HTML}{333333}

\titleformat{\section}{\normalsize\bfseries}{}{0em}{\MakeUppercase}[{\color{rule}\titlerule[0.6pt]}]
\titlespacing*{\section}{0pt}{4.5pt}{2.5pt}

\setlist[itemize]{leftmargin=1.15em, labelsep=0.5em, itemsep=1.4pt, topsep=2pt, parsep=0pt, before=\vspace{1pt}}
\renewcommand{\labelitemi}{\textbullet}

\setlength{\parindent}{0pt}
\setlength{\parskip}{0pt}
\pagestyle{empty}
\linespread{0.96}
\tolerance=1500
\emergencystretch=2.5em

\newcommand{\role}[2]{\textbf{#1}\hfill #2\par}

\begin{document}

% ===== Header =====
\begin{center}
  {\LARGE\bfseries Михайловский Илья Евгеньевич}\\[2pt]
  {\large ML-инженер, NLP и LLM}\\[3pt]
  Москва \textperiodcentered{} офис, гибрид или удалённо \textperiodcentered{} английский C1\\[2pt]
  Telegram: \href{https://t.me/Iluuuaaa}{@Iluuuaaa} \textperiodcentered{}
  \href{mailto:ilyamihailovsy@gmail.com}{ilyamihailovsy@gmail.com} \textperiodcentered{}
  \href{https://github.com/iluuua}{github.com/iluuua}
\end{center}

% ===== Summary =====
\section{О себе}
Обучаюсь на направлении <<Прикладная математика и информатика>> в Московском политехе, выпускаюсь в 2029 году.
Python в продакшене пишу с 2024 года. Основал ReachFlow, платформу ИИ-агентов продаж, и сам реализовал агентов,
LLM-слой диалога и продуктовый бэкенд. Модели писал с нуля на PyTorch. В финале AIDAO 2025 собирал
occupancy-карту в bird's-eye view по камерам беспилотника и получил IoU 0.5606 при 0.564 у победителя олимпиады.
Алгоритмическая база олимпиадная, на C++.

% ===== Skills =====
\section{Навыки}
\begin{itemize}
  \item \textbf{ML и DL.} PyTorch, NumPy, pandas, matplotlib. Обучение с нуля и transfer learning, mixed precision,
        аугментации, подбор лосса под метрику, калибровка порога, дисбаланс классов, ResNet и FPN, сегментация, IoU
  \item \textbf{NLP и LLM.} LLM-пайплайны в проде, prompt engineering, строгий JSON-schema output, классификация
        и extraction, RAG, эмбеддинги и pgvector, агенты и tool calling, MCP, роутинг моделей с фолбэками,
        оценка стоимости и латентности
  \item \textbf{Инженерия.} Python, FastAPI, asyncio, PostgreSQL и SQL, Redis, Docker, CI/CD, Linux, Git, pytest
  \item \textbf{Математика и алгоритмы.} Теория вероятностей, матстатистика, линейная алгебра, методы оптимизации,
        C++, алгоритмы и структуры данных, оценка сложности
\end{itemize}

% ===== Experience =====
\section{Опыт работы}

\role{ReachFlow, ML/NLP и backend инженер}{с января 2026}
\textit{B2B SaaS, ИИ-ассистент, который ведёт продающие переписки в Telegram \textperiodcentered{} \href{https://reachflow.tech}{reachflow.tech}}
\begin{itemize}
  \item Собрал \textbf{конвейер диалога}, где на каждое входящее сообщение уходит восемь вызовов модели.
        Классификатор размечает состояние лида по 16 типам событий с якорем на сообщение и confidence, факты
        извлекаются по 18 типизированным ключам, а ответ пишет стратег, переписывает гуманизатор и проверяет ревьюер.
  \item \textbf{Удешевил ход.} Пять классификаторных вызовов слил в два, потому что одна модель умеет отдать
        несколько секций одним объектом. Вызовов на сообщение стало восемь вместо десяти, а входных символов
        на стадии классификации ушло меньше, от 14 до 43 процентов в зависимости от задачи. Каждая секция
        при этом сохранила свой режим отказа.
  \item Сделал \textbf{строгий структурированный вывод}. В реестре 25 LLM-задач и 22 схемы на pydantic, схемы
        нормализуются под strict-режим вендоров, а ответ проверяется у меня, и его нарушение считается повторяемой ошибкой.
        Исполнитель умеет ретраи и каскад фолбэков между вендорами, каждая попытка попадает в трейс с моделью,
        латентностью и стоимостью.
  \item Собрал \textbf{ИИ-агента с инструментами}. Реестр \texttt{ToolSpec} остаётся единственным источником правды
        сразу для блока инструментов в промпте, гейта привилегий по риску, границы доверия и MCP-эндпоинта.
        У оркестратора есть бюджет вызовов.
  \item Разобрал провал \textbf{семантического гейта}. Косинусная близость на эмбеддингах text-embedding-3-small
        в pgvector с HNSW отбраковывала 55 кандидатов из 63, потому что мерила формат и длину текста вместо
        релевантности. Я перевёл гейт из блокирующего в наблюдательный и оставил решение за LLM-квалификацией.
\end{itemize}
\vspace{2pt}

\role{Фриланс, Python-разработчик}{с июля 2024}
\begin{itemize}
  \item Телеграм-бот-консультант для дистрибьютора пластиковых окон. Внутри \textbf{RAG} на PostgreSQL с векторным
        поиском по базе знаний и интеграция с Битрикс24. По оценке заказчика бот закрыл около 80\% типовых консультаций.
\end{itemize}
\vspace{2pt}

\role{Near AI, разработчик решений на C++}{с июня 2023 по июль 2024}
\begin{itemize}
  \item Готовил алгоритмические решения на C++ как обучающие данные для AI-платформы. Оптимизировал по времени
        и памяти, разбирал асимптотику, ревьюил чужой код в международной команде.
\end{itemize}

% ===== ML projects =====
\section{ML-проекты и соревнования}

\role{AIDAO 2025, международная олимпиада по ИИ и анализу данных (ФКН НИУ ВШЭ, Яндекс Образование)}{ноябрь 2025}
\textit{Финал очно, 32 часа, топ-30 команд из 248 заявок от участников из 14 стран. Задачу ставила команда
беспилотных технологий Яндекса, нужно было построить occupancy-карту в bird's-eye view по камерам автомобиля
со скором по IoU. Решение на PyTorch с нуля.}
\begin{itemize}
  \item Реализовал \textbf{Lift-Splat-Shoot}. Общий для камер ResNet-50 сливает уровни /8 и /16 через FPN,
        depth-head раскладывает глубину на 64 бина с обучаемой температурой softmax, признаки дифференцируемо
        разлетаются в BEV-сетку 188$\times$126. Камеры сливаются через max и mean с SE-блоком, декодер на dilated
        residual-блоках.
  \item Главный прирост дала \textbf{покадровая калибровка} вместо фиксированной. Геометрия проекции пересчитывается
        на каждом кадре из \textit{intrinsics} и \textit{car\_to\_cam}, а сетка пикселей переносится в координаты
        исходного изображения, потому что у четырёх камер разные разрешения.
  \item Лосс собрал из BCE с \texttt{pos\_weight} по балансу классов и $0.3\times$ дифференцируемого surrogate mIoU,
        то есть оптимизировал ровно ту метрику, по которой считался скор. Учил AdamW с раздельными lr для backbone
        и голов, с cosine annealing и AMP, порог калибровал свипом по валидации.
  \item На скрытом тесте вышло \textbf{IoU 0.5606}, у команды-победителя олимпиады 0.564. Инференс 2000 кадров
        занимает 1,7 минуты на одной GPU.
\end{itemize}

% ===== Education =====
\section{Образование и курсы}
\textbf{Московский политехнический университет}, прикладная математика и информатика, выпуск 2029.
\textbf{НИУ ВШЭ, ФКН}, вольнослушатель ПМИ. В 2025 году прошёл \textbf{Deep Learning School} ФПМИ МФТИ,
Январскую научную школу \textbf{ОЦ <<Сириус>>} на 92 академических часа и годовой курс \textbf{Т-Поколения}
<<Алгоритмы>>, параллель B.

\end{document}
